\documentclass[11pt]{article}
\usepackage[margin=1in]{geometry}
\usepackage{amsmath, amssymb, booktabs}
\usepackage[hidelinks]{hyperref}

\title{Residual Trend and Mean-Reversion Extension\\
to the Avellaneda--Lee Daily ETF Replication}
\author{Thiago Amin}
\date{May 2026}

\begin{document}

\maketitle

\section{Motivation}

The baseline replication plan in \texttt{replication\_plan.tex} follows the
Avellaneda--Lee ETF residual strategy: stock returns are decomposed into a
sector ETF component and an idiosyncratic residual, and trades are opened when
the residual level is sufficiently far from its estimated equilibrium. This is
a pure contrarian interpretation of residual displacement. If a stock appears
expensive relative to its sector ETF, the baseline sells the stock and buys the
ETF hedge; if it appears cheap, the baseline buys the stock and sells the ETF
hedge.

That assumption is natural for a paper replication, but it is restrictive in
post-2020 data. The COVID shock, the rate-hiking cycle, and the artificial
intelligence repricing created persistent relative trends inside sectors. In
such environments, a residual that keeps rising may not be temporary noise; it
may reflect a sustained idiosyncratic repricing. A strategy that mechanically
fades every large residual can repeatedly short relative winners too early.

This extension keeps the paper's residual framework but changes the trading
decision. The residual is allowed to be either a mean-reverting displacement
or a relative trend. The strategy first asks which structure is more plausible,
and only then decides whether to fade or ride the residual.

\section{Residual Framework}

For stock $i$ assigned to sector ETF $I(i)$, daily returns are decomposed as
\[
    R_{i,t}
    =
    \beta_{i,t} R_{I(i),t}
    +
    \widetilde{R}_{i,t}.
\]
The coefficient $\beta_{i,t}$ is estimated from a trailing window using only
information available before the trading decision. As in the baseline plan,
the cumulative residual level is
\[
    X_{i,t}
    =
    \sum_{u \le t} \widetilde{R}_{i,u}.
\]
Large positive $X_{i,t}$ means the stock has cumulatively outperformed its
sector ETF after beta adjustment. Large negative $X_{i,t}$ means it has
underperformed.

The baseline strategy assumes that $X_{i,t}$ is stationary and mean-reverting.
This extension treats stationarity as one possible state rather than a
permanent assumption.

\section{Two Residual States}

\subsection{Mean-Reversion State}

In the mean-reversion state, the residual level is interpreted as a temporary
relative-value displacement. The OU score from the baseline plan remains the
decision variable:
\[
    s^{MR}_{i,t}
    =
    \frac{X_{i,t}-\widehat{m}_{i,t}}
         {\widehat{\sigma}_{eq,i,t}}.
\]
A high positive score indicates that the stock is expensive relative to its
sector ETF, so the mean-reversion trade is short stock and long ETF. A high
negative score indicates that the stock is cheap relative to its sector ETF,
so the trade is long stock and short ETF.

\subsection{Trend State}

In the trend state, a persistent residual move is treated as information
rather than as a mispricing to fade. Define a trailing residual trend score
over lookback $L$:
\[
    s^{TR}_{i,t}
    =
    \frac{X_{i,t}-X_{i,t-L}}
         {\widehat{\sigma}_{\Delta X,i,t}\sqrt{L}},
\]
where $\widehat{\sigma}_{\Delta X,i,t}$ is a trailing volatility estimate of
daily residual increments. A positive trend score means the stock is
persistently outperforming its sector ETF. A negative trend score means it is
persistently underperforming.

The trend trade follows the residual:
\[
    s^{TR}_{i,t}>c_{TR}
    \quad\Rightarrow\quad
    \text{long stock, short ETF},
\]
\[
    s^{TR}_{i,t}<-c_{TR}
    \quad\Rightarrow\quad
    \text{short stock, long ETF}.
\]
This is still a relative-value trade. It does not abandon sector hedging; it
rides idiosyncratic momentum against the sector benchmark.

\section{Regime Decision}

The central extension is a mode selection rule. At each decision date, the
strategy chooses between:
\[
    q^{MR}_{i,t}
    =
    -\operatorname{sign}(s^{MR}_{i,t})
    \quad\text{when } |s^{MR}_{i,t}| \text{ is large and trend is weak},
\]
and
\[
    q^{TR}_{i,t}
    =
    \operatorname{sign}(s^{TR}_{i,t})
    \quad\text{when } |s^{TR}_{i,t}| \text{ is large and trend is confirmed}.
\]
Here $q_{i,t}$ is the stock-side direction: $+1$ means long stock and $-1$
means short stock. A simple hard-switching rule is
\[
    q_{i,t}
    =
    \begin{cases}
    \operatorname{sign}(s^{TR}_{i,t}),
        & |s^{TR}_{i,t}| \ge c_{TR},\\
    -\operatorname{sign}(s^{MR}_{i,t}),
        & |s^{TR}_{i,t}| < c_{TR}
          \text{ and } |s^{MR}_{i,t}| \ge c_{MR},\\
    0, & \text{otherwise.}
    \end{cases}
\]
The trend condition is evaluated first. This priority reflects the empirical
problem that motivates the extension: in strong post-2020 repricing regimes,
the costly error is often fading a persistent residual too early.

\section{Volume and Information Confirmation}

The source paper reports that ETF residual strategies improve when volume is
used to distinguish low-volume noise from high-volume information. This idea
fits naturally into the trend-versus-mean-reversion extension.

Let
\[
    \nu_{i,t}
    =
    \frac{V_{i,t}}
         {\overline{V}_{i,t}^{(60)}}
\]
be relative volume. A residual move on unusually high volume is treated as
more likely to contain information. A residual move on ordinary or low volume
is more likely to be a temporary dislocation.

One implementation is to require trend confirmation:
\[
    |s^{TR}_{i,t}| \ge c_{TR}
    \quad\text{and}\quad
    \nu_{i,t} \ge c_V.
\]
Another implementation is to form a trading-time residual increment:
\[
    \widetilde{R}^{TT}_{i,t}
    =
    \frac{\widetilde{R}_{i,t}}{\nu_{i,t}^{\phi}},
    \qquad \phi \ge 0.
\]
This reduces the influence of high-volume residual jumps in the
mean-reversion fit, making the strategy less likely to fade moves that appear
to be information-driven.

\section{Portfolio Construction}

The portfolio construction remains close to the baseline replication. For
stock-side signal $q_{i,t}\in\{-1,0,+1\}$ and base stock allocation
$\Lambda_t$, the stock weight is
\[
    w_{i,t} = \Lambda_t q_{i,t}.
\]
The sector ETF hedge is
\[
    w_{I(j),t}
    =
    -\sum_{i:I(i)=I(j)}
    \widehat{\beta}_{i,t} w_{i,t}.
\]
Thus both residual trend trades and residual mean-reversion trades remain
sector-neutral in the same sense as the paper's ETF strategy. The difference
is the sign of the stock-side trade, not the definition of factor neutrality.

Net daily PNL is still
\[
    \pi_t
    =
    w_t^\top R_t
    -
    c\left\|w_t-w_{t-1}\right\|_1,
\]
with transaction costs charged on position changes.

\section{Diagnostics}

The extension must report performance by trading mode. Aggregate PNL is not
enough, because the key research question is whether residual trend and
residual mean reversion work in different states.

For each date and stock, trades should be labeled as one of:
\[
    \{\text{trend long},\,
      \text{trend short},\,
      \text{mean-reversion long},\,
      \text{mean-reversion short},\,
      \text{flat}\}.
\]
The backtest should then report:
\begin{itemize}
    \item cumulative PNL by mode;
    \item hit rate by mode;
    \item average holding period by mode;
    \item turnover and transaction costs by mode;
    \item sector contribution by mode; and
    \item exposure to the assigned ETF after hedging.
\end{itemize}

If trend-mode PNL dominates, the post-2020 market may be better described by
relative momentum than by residual convergence. If mean-reversion-mode PNL
dominates only in calm regimes, then the strategy should gate contrarian
trades more aggressively. If neither mode works after costs, the residual
signal itself is not strong enough for the chosen universe and horizon.

\section{Implementation Sequence}

This note should be implemented after the baseline assigned-sector-ETF
replication is working. The proposed sequence is:
\begin{enumerate}
    \item keep the assigned sector ETF residual construction from
    \texttt{replication\_plan.tex};
    \item compute residual levels $X_{i,t}$ and OU s-scores as before;
    \item add residual trend scores over one or more fixed lookbacks;
    \item add mode labels for trend, mean reversion, and flat states;
    \item trade with trend priority, then compare against pure mean reversion;
    \item add volume confirmation or trading-time residuals; and
    \item report PNL separately by mode before optimizing thresholds.
\end{enumerate}

The purpose is not to data-mine a better backtest. The purpose is to test a
specific post-2020 hypothesis: some large residuals represent temporary
relative-value dislocations, while others represent persistent idiosyncratic
repricing that should be followed rather than faded.

\end{document}
